\documentclass[aspectratio=169,usenames,dvipsnames]{beamer}
\beamertemplatenavigationsymbolsempty
% \documentclass{beamer}

% UNCOMMENT FOR HANDOUTS
% \usepackage{pgfpages}
% \pgfpagesuselayout{2 on 1}[landscape,letterpaper,border shrink=5mm]
% \setbeameroption{show notes}

\mode<presentation>
\usetheme{Boadilla}
\usecolortheme{spruce}
\usepackage{../../LizStyle,ulem}
\usepackage{color}
\usepackage[color,matrix,arrow]{xy}
% \input xy
% \xyoption{all}
\usepackage{tikz}
\usepackage{tikz-cd}
\tikzset{commutative diagrams/diagrams={ampersand replacement=\&}}

\AtBeginSection{\frame{\sectionpage}}


% ==========Command for putting comments and removing them=======
% Visible:
\newcommand{\InstructorNotes}[1]{\textit{\color{orange}#1}}
\newcommand{\InstructorList}[1]{
\begin{itemize}
	\color{orange}
	#1
\end{itemize}
}
% Removed:
% \renewcommand{\InstructorNotes}[1]{}
% \renewcommand{\InstructorList}[1]{}


\newcommand{\bvfill}{\vskip0pt plus 1filll}

\usepackage{multimedia}
\usepackage{graphbox}

% \usepackage{algpseudocode}


\newcommand{\Reeb}{\textbf{Reeb}}
\newcommand{\Rtopc}{\textbf{Rtopc}}
% \newcommand{\RR}{\mathcal{R}}

\newcommand{\Set}{\textbf{Set}}
\newcommand{\Vect}{\textbf{Vect}}
\newcommand{\Top}{\textbf{Top}}
\newcommand{\Int}{\textbf{Int}}


\graphicspath{ {Figures/} {../Figures/}{../../Figures/} }



\begin{document}
%\pdfoptionpdfminorversion 5
%\logo{\includegraphics[height=1.5cm]{dukeshield.pdf}}
\title[]{Ch 8.2: Ensemble Tree Models}
\subtitle{Lecture 22 - CMSE 381}
\date{Wed, Mar 30, 2022}



\institute[MSU-CMSE]{Michigan State University\\
			::\\
          Dept of Computational Mathematics, Science \& Engineering}
\author[Dr.~Munch]{Prof.~Elizabeth Munch}

\begin{frame}
 \titlepage


\end{frame}
%--------------------------
\begin{frame}[c]{Announcements }

\begin{columns}
\column{.45\textwidth}
\centering

	\textbf{Last time:}
	\begin{itemize}
	\item 8.1 Decision Trees
	\end{itemize}

	\vspace{.2in}

	\textbf{This lecture:}
	\begin{itemize}
	\item
\end{itemize}


\column{.45\textwidth}
\centering

\textbf{Announcements:}
	\begin{itemize}
	\item Mon 3/28: 8.1
	\item Weds 3/30: 8.2
	\item Fri 4/1: Lab: 8.3
	\begin{itemize}
		\item HW \#8 Due Friday
	\end{itemize}

	\end{itemize}
\end{columns}

\end{frame}
%--- Next Frame ---%



%=================
\section{Last time}
%=================




\begin{frame}[c]{First decision tree example}

\begin{columns}
\column{.25\textwidth}
\centering
\includegraphics[width = \textwidth, align = c]{Hitters-TreeSlides-OnlyLog.png}
\column{.35\textwidth}
\footnotesize
\InstructorList{
  \item Top split assigns observations with Years<4.5 to left branch
	\item Return mean response for players with that.
	\item Predictions:
	\InstructorList{
	  \item  mean log salary is 5.107, so returns $\exp(5.107) = \$165.174 $ thousand dollars
		\item $5.999 \Rightarrow \$402,834$
		\item $6.740 \Rightarrow \$845,346$
	}
}

\column{.35\textwidth}
\centering
\includegraphics[width = \textwidth, align = c]{Fig8-1}
\end{columns}

\end{frame}
%--- Next Frame ---%


\begin{frame}[c]{Viewing Regions Defined by Tree}

\begin{columns}
\column{.35\textwidth}
\centering
\includegraphics[width = \textwidth, align = c]{Fig8-1}
\column{.45\textwidth}
\centering

\includegraphics[width = \textwidth, align = c]{Fig8-2}
\end{columns}

\end{frame}
%--- Next Frame ---%


\begin{frame}[c]{How do we actually get the tree? Two steps}

\begin{columns}
\column{.45\textwidth}
\centering
\begin{enumerate}
	\item We divide the predictor space -- that is, the set of possible values for $X_1 , X_2 , \cdots , X _p$ --- into $J$ distinct and non-overlapping regions, $R_1 , R_2 , \cdots, R_J $.
	\item For every observation that falls into the region $R_j$, we make the same prediction $=$ the mean of the response values for the training observations in $R_j $.
\end{enumerate}
\column{.45\textwidth}
\centering
\includegraphics[width = \textwidth, align = c]{Fig8-3b}
\end{columns}

\end{frame}
%--- Next Frame ---%



\begin{frame}[c]{Recursive binary splitting}

\begin{columns}
\column{.3\textwidth}
\centering
\textbf{Goal:}

Find boxes $R_1,\cdots,R_J$ that minimize

\begin{equation*}
	\sum_{j=1}^J \sum_{i \in R_j} (y_i-\hat y_{R_j})^2
\end{equation*}
$\hat y_{R_j} = $ mean response for training observations in $j$th box

\column{.45\textwidth}

Pick $s$ so that splitting into $\{X \mid X_j < s\}$ and $\{X \mid X_j \geq s\}$ results in largest possible reduction in RSS:
\begin{align*}
	R_1(j,s) = \{X \mid X_j < s\}\\
	R_2(j,s) = \{X \mid X_j \geq s\}\\
\end{align*}
\begin{equation*}
	\sum_{i \mid x_i \in R_1(j,s)}(y_i - \hat y _{R_1})^2
	+
	\sum_{i \mid x_i \in R_2(j,s)}(y_i - \hat y _{R_2})^2
\end{equation*}

\column{.2\textwidth}
\includegraphics[width = \textwidth, align = c]{Fig8-3c}
\includegraphics[width = \textwidth, align = c]{Fig8-3b}
\end{columns}

\end{frame}
%--- Next Frame ---%


\begin{frame}[c]{Pruning}

\begin{columns}
\column{.45\textwidth}
\centering
\includegraphics[width = \textwidth, align = c]{Fig8-4}
\column{.45\textwidth}
\centering
\InstructorList{
  \item Big trees leave you open to potential overfitting
	\item Could just stop building earlier, but that's short sighted
	\item Instead, grow a big tree and prune it back
	\item Find subtree with best test error rate
	\item Too many subtrees to test them all
}


\end{columns}

\end{frame}
%--- Next Frame ---%

\begin{frame}[c]{Weakest Link Pruning }
	{Also called Cost complexity pruning}

\begin{columns}
\column{.45\textwidth}
\centering

For every $\alpha$, there is a subtree $T$ that minimizes:
\begin{equation*}
	\sum_{m=1}^{|T|} \sum_{i \mid x_i \in R_m } (y_i - \hat y_{R_m})^2 + \alpha |T|
\end{equation*}

\begin{itemize}
	\item $|T| = $ number of terminal nodes of $T$
	\item $R_m$ is rectangle for $m$th terminal node
	\item $\hat y_{R_m}$ is mean of training observations in $R_m$
\end{itemize}
\column{.45\textwidth}
\centering
\begin{itemize}
  \item $\alpha=0$ gets entire tree
	\item Increasing $\alpha$ penalizes size of tree
	\item Branches pruned from tree in nested and predictable fashion
	\item Easy to get trees for all vales of $\alpha$
	\item Pick $\alpha $ via CV
\end{itemize}

\includegraphics[width = .49\textwidth, align = c]{Fig8-4}
\includegraphics[width = .49\textwidth, align = c]{Fig8-1}
\end{columns}

\end{frame}
%--- Next Frame ---%
\begin{frame}[c]{Algorithm version}

\begin{columns}
\column{.65\textwidth}
\centering
\includegraphics[width = \textwidth, align = c]{Alg8-1}

\end{columns}

\end{frame}
%--- Next Frame ---%
\begin{frame}[t]{Classification version}

\begin{columns}
\column{.45\textwidth}
\centering

\includegraphics[width = \textwidth, align = c]{Fig8-6a}

\includegraphics[width = .49\textwidth, align = c]{Fig8-6b}
\includegraphics[width = .49\textwidth, align = c]{Fig8-6c}
\column{.45\textwidth}
\centering
Evaluating the splits:
\begin{itemize}
	\item $\hat p_{mk} = $ proportion of training observations in $R_m$ from the $k$th class
	\item Error: $E = 1-\max_{k}(\hat p_{mk})$
	\item Gini index:
	\begin{equation*}
		G = \sum_{k=1}^K \hat p_{mk} (1-\hat p_{mk})
	\end{equation*}
	\item Entropy:
	\begin{equation*}
		D = - \sum_{k=1}^K \hat p_{mk} \log \hat p _{mk}
	\end{equation*}
\end{itemize}

\end{columns}

\end{frame}
%--- Next Frame ---%

%=====
\section{Ensemble Tree Methods}
%=====







\begin{frame}[c]{TL:DR}


\end{frame}
%--- Next Frame ---%




\begin{frame}[t]{Next time}


\begin{columns}
\column{.45\textwidth}
\centering

\textbf{Next time:}

\begin{itemize}
	\item Lab: Trees
\end{itemize}

\column{.45\textwidth}
\centering

\textbf{This week:}
\begin{itemize}

	\item Mon 3/28: 8.1
	\item Weds 3/30: 8.2
	\item Fri 4/1: Lab: 8.3
	\begin{itemize}
		\item HW \#8 Due Friday
	\end{itemize}
\end{itemize}
	\end{columns}

\end{frame}
%--- Next Frame ---%


\end{document}
